%=====================================================================
\section{Introduction}
\label{sec:intro}

%---------------------------------------------------------------------
\begin{figure*}[t]
\centering
\resizebox{\textwidth}{!}{%
\begin{tikzpicture}[
    font=\small,
    >={Stealth[length=2mm]},
    rule/.style={draw, rounded corners=2pt, fill=gray!15, minimum height=5.5mm, inner sep=2.5pt, font=\footnotesize\ttfamily},
    dec/.style={draw, rounded corners=2pt, fill=blue!10, minimum height=5.5mm, inner sep=2.5pt, font=\footnotesize\ttfamily},
    box/.style={draw, rounded corners=3pt, inner sep=5pt, align=left},
    ttl/.style={font=\small\bfseries, anchor=south west},
  ]
  % ---------------- (a) the procedure
  \node[dec] (st) at (2.6,0) {sent\_type};
  \node[dec] (pr) at (0.7,-1.2) {predicate};
  \node[dec] (np) at (2.6,-1.2) {np\_phrase};
  \node[dec] (se) at (4.65,-1.2) {special\_entity};
  \node[rule] (cov) at (0.7,-2.3) {coverage};
  \node[rule] (ant) at (2.6,-2.3) {antecedent};
  \draw[->] (st) -- (pr); \draw[->] (st) -- (np); \draw[->] (st) -- (se);
  \draw[->] (pr) -- (cov); \draw[->] (pr) -- (ant);
  \draw[->, dashed] (pr.north) to[bend left=35] node[font=\scriptsize\itshape, above left, pos=0.4] {open slots} (st.west);
  \node[rule] (re) at (0.7,-3.5) {reentrancy};
  \node[dec]  (doc) at (2.6,-3.5) {document};
  \node[rule] (out) at (4.65,-3.5) {assemble};
  \draw[->] (re) -- (doc); \draw[->] (doc) -- (out);
  \node[font=\scriptsize\itshape, anchor=west] at (-0.4,-2.9) {later waves};
  \draw[gray, densely dotted] (-0.4,-2.75) -- (5.9,-2.75);
  \node[rule, minimum height=3mm, minimum width=4mm] (lr) at (-0.2,-4.35) {};
  \node[font=\scriptsize, anchor=west] at (lr.east) {rule node (code)};
  \node[dec, minimum height=3mm, minimum width=4mm] (ld) at (2.5,-4.35) {};
  \node[font=\scriptsize, anchor=west] at (ld.east) {decision node (model)};
  \node[ttl] at (-0.5,0.45) {(a) The guideline as a procedure};
  % ---------------- (b) a closed decision
  \node[box, fill=blue!4, text width=6.6cm, anchor=north west] (q) at (6.6,0.3) {%
    \textbf{Code enumerates} the content words of the clause and the PropBank senses of each verb.\\[3pt]
    \textbf{Q1} Which word heads the clause?\\
    \quad \emph{explosion} .07 \quad \emph{killed} \textbf{.91} \quad \emph{people} .02\\[1pt]
    \textbf{Q2} Sense of \emph{kill}? \quad kill-01 \textbf{.88} \quad kill-02 .09\\[1pt]
    \textbf{Q3} Role of \emph{explosion}? \quad ARG0 \textbf{.71} \quad none .18\\[3pt]
    \textbf{Rules build} \texttt{(k / kill-01 :ARG0 ...)}, reading each answer against a threshold $\theta$. An answer with confidence $<\tau$ can be routed to a stronger model or to a person.};
  \node[ttl] at ([yshift=2pt]q.north west) {(b) A closed decision inside \texttt{predicate}};
  % ---------------- (c) learning the procedure
  \node[box, fill=green!6, text width=4.8cm, anchor=north west] (l) at (13.9,0.3) {%
    \textbf{Given} 1--4 gold documents:\\[3pt]
    \textbf{Calibrate} the thresholds $\theta$ by re-parsing and scoring.\\[2pt]
    \textbf{Route} decisions with confidence $<\tau$.\\[2pt]
    \textbf{Reflect}: find each difference from gold, attribute it to a node through the trace, edit that node (text, options, rules, $\theta$, registry), re-parse, and keep the edit only if the score rises.\\[3pt]
    \textbf{Result}: a new version of the procedure, as a readable diff.};
  \node[ttl] at ([yshift=2pt]l.north west) {(c) Learning the procedure};
  % ---------------- arrows between panels
  \draw[->, thick] (se.east) -- ++(0.5,0) |- ([yshift=4mm]q.west);
  \draw[->, thick] (q.east |- l.west) -- (l.west);
  \coordinate (low) at ([yshift=-3mm]l.south);
  \draw[->, thick, dashed] (l.south) -- (low) -| node[font=\scriptsize\itshape, pos=0.3, below] {edits: a new version of the procedure} (out.south);
\end{tikzpicture}}
\caption{Decision orchestration. (a)~The guideline is compiled into skills dispatched by a registry. Rule nodes are code; decision nodes consult a model. Shown is part of the \umr procedure (all nodes: Table~\ref{tab:nodes}, appendix). (b)~In closed mode, a decision node asks questions whose options code has enumerated, and rules read the answer probabilities. (c)~With a handful of gold documents, the procedure itself is learned. Example sentence, in English for readability: \emph{Police said the explosion killed two people.}}
\label{fig:overview}
\end{figure*}

New annotation schemes appear all the time: an intent taxonomy for a product, an event schema for a news domain, a coding scheme in political science, a meaning representation for an under-documented language. Each starts with a guideline and almost no annotated data, and producing the first data is the bottleneck. Annotators must learn the guideline and label item after item before a model can be trained to take over \citep{settles2009active}. For rich semantic representations the cost is extreme. Uniform Meaning Representation \citep[\umr;][]{gysel2021design} requires a full graph for every sentence and a document-level layer of temporal, modal and coreference relations, and after years of effort its release has only a few hundred gold graphs per major language \citep{bonn2025umr2}.

LLMs promise to shortcut this cold start, and for labels from a familiar inventory they often match crowd workers \citep{wang2021reduce,gilardi2023chatgpt,tan2024annotation}. Structured annotation is harder. Asked for a whole \umr or \amr graph, LLMs ignore the conventions the guideline defines, produce malformed graphs, and entangle their errors so that no one can tell which part of the guideline failed \citep{ettinger2023expert,sun2024chinese}. In our experiments a frontier model given the sentence and the guideline reaches about 30 sentence \F, and an open 9B model produces no graph that can be scored at all. Training on annotated data is what the cold start lacks. Agent harnesses that let the model plan its own steps \citep{yao2023react} leave cost and output format to the model, hide which decision caused an error, and work only with the strongest models.

Our starting point is that \emph{a guideline is mostly a procedure}. To annotate a sentence in \umr, one segments it into clauses, finds each predicate, looks up its frame, attaches arguments, writes names and dates in fixed patterns, merges mentions of one referent, keeps the graph well-formed and every word covered, and finally relates each event to time, to its source and to earlier mentions. Most of these steps can be executed by code. A few are judgments that no rule captures: is this word an event here, which sense is meant, is the aspect a state or a performance, does this mention refer to that one?

We propose \textbf{decision orchestration} (Figure~\ref{fig:overview}). The guideline is compiled into a registry of \emph{skills} that a deterministic harness executes. \emph{Rule nodes} are code. \emph{Decision nodes} consult a model, and the model only decides. It either fills a schema-constrained fragment that the harness assembles, or it answers \emph{closed questions} (yes/no, or a choice among candidates that code enumerated) with a probability per option, which rules turn into structure. The model never chooses the next step or writes the final format, and every decision is recorded with its input. The same interface can be served by a frontier LLM, a small open model, or a \emph{decision model}, a service that answers closed questions with probabilities and generates no text \citep{typesafe2026jev,openai2026decisions,microsoft2026decision}.

The procedure, not a model, then becomes the thing to learn. Because it is explicit and every decision is attributable, a handful of gold documents is enough to improve it, without updating any weights (\S\ref{sec:learning}). \textbf{Calibration} sets the thresholds through which rules read decision probabilities by measurement on a few documents, which is needed whenever the decision model changes. \textbf{Routing} sends the few uncertain decisions to a stronger model or, in a cold start, to a person. \textbf{Reflection} has an LLM compare the output with gold, follow the execution trace to the node that caused each difference, and edit that node: its instructions, candidate options, assembly rules, thresholds, or the registry. When parsing is cheap, each edit is a controlled trial on training, regression and development documents, kept only if it helps. Every learned change is a readable diff.

On Chinese \umr parsing (sentence graphs and document-level relations), learning from one to four documents (51--201 sentences), we find that (i) orchestration makes structured annotation feasible for small models: Qwen3.5-9B goes from 0 to 17.6 \F, a 9B decision model reaches 30.8, and a 27B open model reaches 94\% of a frontier model; (ii) the procedure learns from very little data: two gold documents raise a frontier model from 45.3 to 52.6 \F, and the same loop improves a decision-model procedure on documents it never saw; (iii) decisions can be calibrated: recalibrating on four documents gives a new decision model +2.9 \F, and routing low-confidence decisions recovers 75\% of a frontier model's gain at 26\% of its cost; and (iv) errors stay debuggable: lost credit is attributed to skills, measured reflection rejects 40\% of its own proposals, and swapping one component at a time locates the bottleneck. We also show where the approach stops paying off, since gains on unseen documents are an order of magnitude smaller than on training documents. \todo{one sentence on the second benchmark (E9)}

Our contributions are \textbf{decision orchestration}, a recipe for compiling a guideline into a procedure in which code owns control flow and well-formedness and the model only decides; \textbf{a way to learn the procedure from a handful of gold documents} by calibration, routing, and attributed, measured reflection, with no weight updates; and \textbf{an empirical study} across open and frontier LLMs and three decision-model services of what the procedure learns, what transfers, and what it costs.
